\subsection{Neutrino Masses and Oscillations in Models with Large Extra Dimensions --- arXiv:hep-ph/9907520}

\textbf{Authors:} R. N. Mohapatra, S. Nandi, A. P\'erez-Lorenzana

\paragraph{主な主張}
brane上の右巻きニュートリノとbulk singletを含む左右対称模型で、質量ゼロの軽い固有状態とKK混合によるニュートリノ消失を構成し、結合と対称性の破れのスケールに関する仮定の下で基本スケール$M_*\gtrsim10^8\,\mathrm{GeV}$を導く\cite{Mohapatra:1999BulkOscillations}。

\paragraph{新規性}
active neutrinoを直接bulkに結合する最小設定とは異なり、brane上の右巻きニュートリノを介して、軽い状態とKK塔の質量・混合の構造を変える点にある。

\paragraph{位置付け}
bulk neutrino模型の初期の代替構成であり、余剰次元によるニュートリノ質量生成を単一の最小Dirac模型と同一視しないための比較文献である。
